\documentclass[10pt,a4paper]{article}
\usepackage{fontspec,polyglossia}
\setmainlanguage{russian}
\setotherlanguage{english}
\setmainfont{Times New Roman}
\setmonofont{DejaVu Sans Mono}
\usepackage{amsmath,amssymb,booktabs,tabularx,graphicx}
\usepackage[margin=19mm]{geometry}
\usepackage[hidelinks]{hyperref}
\setlength{\parindent}{0pt}
\setlength{\parskip}{4pt}
\setlength{\emergencystretch}{3em}
\author{}\date{}
\title{Обученная синхронизация сети связанных осцилляторов}
\begin{document}\maketitle

\section*{Цель эксперимента}

Нужна динамическая система, в которой обучение приводит к упрощению поведения, но абсолютное число активных компонент заранее неизвестно. Мы взяли сеть из 32 неоднородных связанных осцилляторов. До обучения узлы движутся несогласованно, а после обучения переходят к коллективному режиму.

Главный вопрос: \textbf{может ли scalar-log протокол обнаружить этот переход по одному скалярному датчику?} Полное состояние сети используется только для независимой проверки.

\section*{Что обучается}

У осциллятора $i$ есть фаза $\theta_i(t)$ и собственная частота $\omega_i$. Фиксированный циклический граф задаёт матрицу связей $A$. Динамика имеет вид

\[
\dot\theta_i=\omega_i+(0.10+g)\sum_j A_{ij}\sin(\theta_j-\theta_i).
\]

Здесь $g\geq 0$ --- единственный обучаемый параметр: дополнительная сила связи между узлами. Частоты неоднородны: $\omega_i=1+0.20\,\varepsilon_i$, где $\varepsilon_i$ фиксированы для каждого запуска.

На каждом шаге обучения запускаем десять траекторий с разными начальными фазами. Используем порядок синхронизации

\[
R(t)=\left|\frac{1}{N}\sum_i e^{\mathrm{i}\theta_i(t)}\right|
\]

и минимизируем

\[
L(g)=\frac{1}{10}\sum_{m=1}^{10}(1-R_m)^2+0.006g^2.
\]

Таким образом, обучение решает конкретную задачу управления: \textbf{сделать движение сети согласованным при ограниченной силе связи}. Наш метод не участвует в обучении и не использует $R(t)$ для оценки размерности.

Обучение повторено для трёх seeds. Полученные значения $g$ равны 1.840, 1.746 и 1.814.

\section*{Как применяется scalar-log протокол}

Для каждого seed используем один и тот же фиксированный агрегированный датчик до и после обучения:

\[
s(t)=\frac{1}{\sqrt{N}}\left(a^\top\sin\theta(t)+b^\top\cos\theta(t)\right),
\]

где $a$ и $b$ выбираются один раз и не меняются между этапами. Оценщик не видит фазы отдельных узлов и не получает порядок $R(t)$.

Рассматриваем три этапа:

\begin{center}\small
\begin{tabularx}{\linewidth}{>{\raggedright\arraybackslash}X>{\raggedright\arraybackslash}X>{\raggedright\arraybackslash}X}
\toprule
Этап & Коэффициент & Смысл \\
\midrule
до обучения & $g=0$ & исходный несинхронный режим \\
середина & $g=0.45g_{\mathrm{final}}$ & частичная синхронизация \\
после обучения & $g=g_{\mathrm{final}}$ & результат обучения \\
\bottomrule\end{tabularx}\end{center}

Систему моделируем 18 000 шагов, первые 2 000 отбрасываем. Для всех этапов фиксируем одну конфигурацию оценщика: $E_{\max}=12$, $\tau=20$, $k=8$, исключение Тейлера 700 и окно $W=8192$.

\section*{Результаты}

\begin{center}
\includegraphics[width=\linewidth]{overview.pdf}
\end{center}

Ниже приведены медианы по трём seeds.

\begin{center}\small
\begin{tabularx}{\linewidth}{>{\raggedright\arraybackslash}X>{\raggedright\arraybackslash}X>{\raggedright\arraybackslash}X>{\raggedright\arraybackslash}X>{\raggedright\arraybackslash}X>{\raggedright\arraybackslash}X}
\toprule
Этап & Средний $R$ & SD$(R)$ & PR полного состояния & MG по одному скаляру & LB по одному скаляру \\
\midrule
до обучения & 0.225 & 0.108 & 28.1 & 13.71 & 34.71 \\
середина & 0.835 & 0.016 & 2.02 & 1.30 & 1.35 \\
после обучения & 0.226 & 0.00001 & 2.00 & 1.30 & 1.35 \\
\bottomrule\end{tabularx}\end{center}

У третьего seed после обучения средний $R$ оказался ниже из-за выбора фазового кластера, но его стандартное отклонение практически исчезло. Поэтому все три запуска перешли к почти периодическому режиму, хотя абсолютная фаза и кластер различались.

Переход подтверждается тремя независимыми признаками:

1. порядок $R(t)$ становится почти постоянным;
2. effective rank полного состояния падает примерно с 20--30 до 2;
3. scalar-log оценка MG падает примерно с 11--14 до 1.1--1.3.

Оценка LB показывает тот же переход. Окна не помечаются как вырожденные. Падение поэтому не объясняется только отказом реконструкции или отсутствием соседей.

\section*{Интерпретация}

Мы не утверждаем, что $MG$ равен PR полного состояния. Это разные величины. PR использован как независимое свидетельство того, что движение всей сети стало проще, а scalar-log оценка отвечает на другой вопрос: виден ли этот переход из одного наблюдаемого сигнала.

В этой постановке метод обнаружил \textbf{момент упрощения} без знания абсолютного числа компонент. До обучения сеть занимала многомерный режим, после обучения её движение стало почти синхронным и периодическим. Тот же переход виден в одном скалярном сигнале.

\section*{Ограничения и следующий шаг}

Это контролируемый proof of concept. Полное состояние здесь маленькое, поэтому вычислительное преимущество ещё не измерено. В большой распределённой системе можно постоянно хранить один или несколько агрегированных датчиков, а дорогой анализ всех состояний, спектра Якобиана или показателей Ляпунова запускать только на выбранных этапах.

В текущем эксперименте обучался глобальный коэффициент связи, одинаковый для всех узлов, а граф был задан заранее. Следующая версия должна использовать локальные обучаемые регуляторы, большие сети, шум и пропуски датчиков. Отдельно нужно сравнить время и память scalar-log анализа с полным анализом состояния.

Итоговый проверенный вывод: \textbf{в обучаемой рекуррентной системе синхронизация, подтверждённая полным состоянием, сопровождается резким падением scalar-log оценки по одному агрегированному сигналу.}
\end{document}